% Trust Index paper -- LaTeX source for arXiv submission.
%
% Converted from compliance/TRUST_INDEX_PAPER.md (Markdown draft) on
% 2026-08-14. Content is a direct translation, not a rewrite -- if the
% two ever disagree, treat the Markdown draft as having drifted and
% re-sync it from here (or vice versa), rather than trusting either
% blindly. See TRUST_INDEX_PAPER.md's own header for what still needs
% a human before this can actually be submitted (endorser, independent
% reader) -- this file only satisfies the "convert to LaTeX" step.
%
% Uses the plain `article` class deliberately: arXiv accepts it across
% every plausible category for this paper (cs.AI, cs.CY), and it avoids
% any dependency on a conference/journal class that might not be
% installed in arXiv's compilation environment. Re-verify current arXiv
% format requirements before uploading -- these can change.

\documentclass[11pt]{article}

\usepackage[margin=1in]{geometry}
\usepackage{times}
\usepackage{amsmath}
\usepackage{hyperref}
\usepackage{enumitem}
\usepackage{booktabs}
\usepackage{microtype}
\usepackage[hyphens]{xurl}

\hypersetup{
  colorlinks=true,
  linkcolor=black,
  citecolor=black,
  urlcolor=blue,
}

% \texttt{} cannot break at `/` or `_`, which overflows the margin on
% long file paths (see Section 4). \code{} reuses \url{}'s line-breaking
% behavior, styled monospace via \urlstyle{tt} below, without treating
% the content as a clickable hyperlink.
\urlstyle{tt}
\newcommand{\code}[1]{\url{#1}}

\title{The Trust Index: An Open, Citable Methodology for\\
Multi-Dimensional AI System Trustworthiness Scoring}

\author{WhitePact\thanks{\url{https://github.com/Guruprasath-Annadurai/Whitepact}}}

\date{}

\begin{document}

\maketitle

\begin{abstract}
AI governance today lacks a standardized, openly-verifiable signal for
comparing the trustworthiness of deployed AI systems across
organizations --- unlike domains such as payment security (PCI-DSS) or
web application security (OWASP Top 10), where a published, versioned
standard lets any party independently assess and cite compliance
without depending on a single vendor's proprietary methodology. We
present the \textbf{Trust Index}, an open specification for scoring an
AI system across six weighted dimensions --- fairness, privacy,
security, robustness, compliance, and authenticity --- producing a
0--100 composite score, a letter grade, and a risk tier. We describe
three distinct provenance paths for a Trust Index score
(self-assessment, automated measurement against a public inference
API, and human-reviewed certification), each explicitly labeled to
prevent the conflation of self-reported and independently-verified
claims --- a distinction we argue is necessary for any trust signal to
remain meaningful as adoption scales. We provide a reference
implementation (\texttt{src/responsibleai/trust/score.py}) and a
public verification mechanism
(\texttt{GET /api/trust-index/verify/\{id\}}) that returns the full
scored record for any cited passport ID, making every citation
independently checkable rather than an unfalsifiable claim. We discuss
the specification's current limitations, including the absence of a
formal third-party accreditation process and the lack of statistical
confidence intervals on point-estimate scores, and outline a
versioning discipline intended to keep the published standard
synchronized with its reference implementation over time.
\end{abstract}

\section{Introduction}

As organizations increasingly deploy large language models and other
AI systems in production, buyers, regulators, and end users face a
recurring question with no standardized answer: \emph{how trustworthy
is this specific AI system, on what basis, and who verified the
claim?} Existing approaches tend toward one of two extremes. At one
end, informal marketing claims (``our AI is safe and unbiased'') carry
no verifiable structure at all. At the other, formal regulatory
compliance frameworks (the EU AI Act, NIST AI RMF) are comprehensive
but not designed to produce a single, comparable, publicly citable
score suitable for quick buyer-side comparison the way a credit rating
or a PCI-DSS attestation is.

We propose the Trust Index to fill this gap: a lightweight, open,
versioned scoring methodology, deliberately modeled on the publication
pattern of security standards bodies (OWASP, PCI Security Standards
Council) rather than on a proprietary vendor certification. The
specification is published in full
(\texttt{compliance/TRUST\_INDEX\_SPEC.md}) and its reference
implementation ships as open-source code
(\texttt{src/responsibleai/trust/score.py}), so that the formula
computing a published score is never separated from a paper
description that could drift from what is actually run.

This paper's contribution is threefold: (1) a six-dimension weighted
scoring methodology with defined grade and risk-tier bands; (2) an
explicit three-tier provenance model distinguishing self-reported,
automatically measured, and human-certified scores, which we argue is
necessary to prevent the standard's erosion into an unfalsifiable
marketing claim as adoption grows; and (3) a public verification
mechanism that makes any cited score independently checkable by any
third party, without requiring trust in the citing party alone.

\section{Related work}

\textbf{Regulatory AI risk frameworks.} The EU AI Act \cite{euaiact}
and the NIST AI Risk Management Framework \cite{nistairmf} provide
comprehensive governance requirements but are structured as compliance
obligations and voluntary guidance respectively, not as a single
comparable numeric score intended for rapid, informal buyer-side
comparison across vendors.

\textbf{Model evaluation benchmarks.} Benchmarks such as TruthfulQA
\cite{truthfulqa}, BBQ \cite{bbq}, and HellaSwag \cite{hellaswag}
measure specific narrow capabilities (truthfulness, bias in
question-answering, commonsense reasoning) and are widely used in the
research community, but are not packaged as a composite,
governance-oriented trust score, nor do they distinguish self-reported
from independently measured results as a first-class feature of the
standard itself.

\textbf{Industry security standards.} PCI-DSS \cite{pcidss} and
ISO/IEC~27001 \cite{iso27001} demonstrate the publication pattern this
work follows --- an open, versioned standard separable from any single
accreditation body, with certification as an optional paid layer atop
a freely available specification. We adopt this structure directly,
substituting AI trustworthiness dimensions for payment or
information-security controls.

\textbf{This work's distinguishing claim} is not a novel scoring
formula in isolation, but the combination of (a) an open, versioned,
code-synchronized specification, (b) an explicit, labeled three-tier
provenance model, and (c) a public, per-score verification endpoint
--- together intended to keep a ``Trust Index score'' a checkable
claim rather than a marketing assertion, even as the standard is cited
by parties other than its original publisher.

\section{Methodology}

\subsection{Dimensions and weighting}

A Trust Index score is a weighted composite across six dimensions,
each normalized to $[0, 1]$ before weighting (Table~\ref{tab:dims}).

\begin{table}[h]
\centering
\caption{Trust Index dimensions and default weights}
\label{tab:dims}
\begin{tabular}{@{}llp{8cm}@{}}
\toprule
Dimension & Weight & Definition \\
\midrule
Fairness & 0.20 & Absence of detected bias across protected categories in system outputs. \\
Privacy & 0.15 & Avoidance of leaking or mishandling personally identifiable information. \\
Security & 0.20 & Resistance to adversarial manipulation (prompt injection, jailbreaks, data exfiltration, role-confusion). \\
Robustness & 0.15 & Factual reliability; resistance to hallucination. \\
Compliance & 0.20 & Regulatory/governance process maturity (documented audit trails, incident response, framework alignment). \\
Authenticity & 0.10 & For media-generating/evaluating systems: resistance to deepfake/synthetic-media misuse. Not applicable to text-only systems. \\
\bottomrule
\end{tabular}
\end{table}

The overall score is computed as:
\begin{equation}
\text{overall\_score} = \left(\sum_{i=1}^{6} \text{dimension\_value}_i \times \text{weight}_i\right) \times 100
\end{equation}
producing a value in $[0, 100]$. Grade bands are defined as A $\geq$
90, B $\geq$ 80, C $\geq$ 70, D $\geq$ 60, F $<$ 60; risk tiers as LOW
$\geq$ 80, MEDIUM $\geq$ 60, HIGH $\geq$ 40, CRITICAL $<$ 40.

\subsection{Handling not-applicable dimensions}

Not every system under evaluation admits a meaningful measurement on
every dimension --- a text-only system has no media-authenticity
signal, for instance. Rather than silently omitting a dimension (which
would alter the effective weighting distribution non-transparently) or
fabricating a measurement, we specify that a not-applicable dimension
be held at a disclosed neutral value of 0.5, with the disclosure
itself a required part of any citation. This design choice trades a
small loss of precision for transparency about what was actually
measured versus assumed neutral.

\subsection{Provenance: three distinct paths to a score}

We argue that the central design risk for any open, citable trust
standard is \emph{provenance conflation} --- the failure to
distinguish a self-reported claim from an independently verified one
once both produce numbers in the same visible format. We address this
directly by defining three structurally distinct, separately labeled
paths to a Trust Index score:

\begin{enumerate}
\item \textbf{Self-assessment} (free, unverified): a party submits its
  own dimension values and receives a permanently recorded, hashed
  record. The arithmetic is verifiable; the underlying input claims
  are not audited. Every self-assessed record is labeled
  \texttt{certified: false} everywhere it is displayed or returned via
  API.
\item \textbf{Automated measurement} (free, independently measured,
  narrower scope): dimensions are computed by directly querying a
  system's public inference API against a fixed evaluation corpus,
  without relying on self-reported values for the dimensions it
  covers. This is more credible than self-assessment but is only
  applicable to systems reachable through a public API, and --- in our
  reference implementation --- still holds the compliance and
  authenticity dimensions at the disclosed neutral placeholder, since
  those are not directly measurable via API probing alone.
\item \textbf{Human-reviewed certification} (paid, no automated path):
  a reviewer examines the evidence behind a submitted score and, if it
  holds up, marks the record \texttt{certified: true} with a named
  certifier and timestamp. We deliberately provide no automated route
  to this state --- the moment certification could be self-served, it
  would cease to carry information beyond self-assessment.
\end{enumerate}

\subsection{Verification mechanism}

Every generated record receives a stable identifier and a SHA-256
verification hash. A public endpoint
(\texttt{GET /api/trust-index/verify/\{passport\_id\}}) returns the
complete stored record for any identifier; an unresolvable identifier
returns an explicit 404, signaling that a citation referencing it is
unverifiable. We treat this mechanism --- not the scoring formula
alone --- as the paper's central methodological contribution: a trust
claim that cannot be checked by a party other than the one making it
provides no more assurance than an unstructured marketing statement,
regardless of how principled its underlying arithmetic is.

We additionally define a minimal citation format (Section~3.5)
intended to make every citation self-documenting with respect to
provenance, and an embeddable badge mechanism
(\texttt{GET /api/trust-index/badge/\{id\}.svg}) that renders the same
provenance distinction visually (``Self-Assessed'' vs.\ ``Certified'')
wherever a score is displayed outside the verification page itself.

\subsection{Citation format}

We specify a minimum citation format requiring the scored value,
grade, specification version, explicit provenance label, and a
verification URL:

\begin{quote}
\emph{``[System] scored [X]/100 (Grade [Y]) under the [Standard Name]
v[version] --- [self-reported | certified by [certifier] on [date]].
Verify at [verify\_url].''}
\end{quote}

\section{Reference implementation}

The methodology described here is implemented as open-source code.
Scoring and provenance labeling are implemented across four modules,
one per concern rather than one monolith:
\code{src/responsibleai/trust/score.py} computes the weighted
composite (Section~3.1) from six dimension values;
\code{src/responsibleai/trust/passport.py} generates the immutable,
hashed record for a computed score;
\code{src/responsibleai/db/passport_repository.py} persists records
and serves the verification lookup described in Section~3.4; and
\code{src/responsibleai/leaderboard/runner.py} implements the
automated-measurement provenance path of Section~3.3 specifically ---
querying a model's public inference API against the TruthfulQA corpus
to compute the robustness dimension live, while holding compliance and
authenticity at the disclosed neutral placeholder described in
Section~3.2, since neither is directly observable from API probing
alone (see \code{compliance/LEADERBOARD_METHODOLOGY.md} for the
full measurement protocol). The three provenance paths therefore
correspond to three distinct, identifiable code paths rather than a
single generator parameterized after the fact --- self-assessment and
certification are served by \code{POST /api/trust-index/assess} and
\code{POST /api/trust-index/certify/{id}} respectively, and
automated measurement by \code{GET /api/leaderboard} and its
underlying \texttt{LeaderboardRunner}.

We consider synchronization between the published specification and
the running implementation a first-class design requirement, not an
afterthought: the specification document states explicitly that the
reference implementation's default weights are the source of truth
for the current version, and specification version bumps are recorded
as dated changelog entries alongside the corresponding code change,
rather than as an independently drifting paper standard.

\section{Discussion and limitations}

\textbf{No formal third-party accreditation process yet.} As of the
current specification version, certification is performed only by the
standard's originating team. We consider this an honest, stated
limitation rather than a design goal --- a mature version of this
standard would define how an independent auditor could become an
accredited certifier, analogous to accreditation bodies for ISO
management-system standards, but this is not yet built.

\textbf{Point estimates without confidence intervals.} A Trust Index
score is a point estimate from a single evaluation run, not a
distribution with error bars. Future versions may incorporate
repeated-measurement variance, particularly for the
automated-measurement provenance path, where re-querying a live
model's API could in principle surface score volatility that a single
measurement does not capture.

\textbf{Not a substitute for regulatory compliance certification.}
Scoring well under this standard is not itself a compliance
certification for any specific law or regulation. The compliance
dimension references frameworks --- specifically the NIST AI Risk
Management Framework, the EU AI Act, and ISO/IEC~42001
(\texttt{src/responsibleai/compliance/engine.py}) --- as inputs to its
own scoring, but a high Trust Index score should not be represented,
by any party citing it, as equivalent to formal regulatory compliance.

\textbf{Self-assessment remains gameable in principle.} Nothing
prevents a self-assessing party from submitting favorable, unverified
dimension values. We view this as an accepted, disclosed limitation of
the self-assessment tier specifically (matching the explicit
\texttt{certified: false} labeling), not a flaw in the standard as a
whole --- the same way a self-reported SAT score is arithmetically
real but not independently proctored, and is understood as such by
anyone reading it correctly labeled.

\section{Conclusion}

We present the Trust Index as an open, versioned, code-synchronized
methodology for scoring AI system trustworthiness, distinguished from
prior work primarily by its explicit three-tier provenance model and
public per-score verification mechanism. We publish the full
specification and reference implementation openly, in the pattern of
established security and payment-industry standards, with the goal
that a Trust Index citation become a checkable claim rather than a
marketing assertion, regardless of which party is doing the citing. We
identify the absence of a formal third-party accreditation process as
the standard's most significant current limitation and an explicit
direction for future work.

\begin{thebibliography}{9}

\bibitem{truthfulqa}
Lin, S., Hilton, J., \& Evans, O. (2022).
TruthfulQA: Measuring How Models Mimic Human Falsehoods.
In \emph{Proceedings of the 60th Annual Meeting of the Association for
Computational Linguistics (Volume 1: Long Papers)}, pages 3214--3252,
Dublin, Ireland. Association for Computational Linguistics.
\url{https://doi.org/10.18653/v1/2022.acl-long.229}
(arXiv:2109.07958)

\bibitem{bbq}
Parrish, A., Chen, A., Nangia, N., Padmakumar, V., Phang, J.,
Thompson, J., Htut, P.~M., \& Bowman, S.~R. (2022).
BBQ: A Hand-Built Bias Benchmark for Question Answering.
In \emph{Findings of the Association for Computational Linguistics:
ACL 2022}, pages 2086--2105, Dublin, Ireland. Association for
Computational Linguistics. (arXiv:2110.08193)

\bibitem{nistairmf}
Tabassi, E. (2023).
Artificial Intelligence Risk Management Framework (AI RMF 1.0).
NIST Trustworthy and Responsible AI, NIST AI 100-1, National Institute
of Standards and Technology, Gaithersburg, MD.
\url{https://doi.org/10.6028/NIST.AI.100-1}

\bibitem{euaiact}
European Parliament and Council of the European Union. (2024).
Regulation (EU) 2024/1689 of the European Parliament and of the
Council of 13 June 2024 laying down harmonised rules on artificial
intelligence (Artificial Intelligence Act).
\emph{Official Journal of the European Union}, published 12 July 2024.
\url{https://eur-lex.europa.eu/eli/reg/2024/1689/oj}

\bibitem{pcidss}
PCI Security Standards Council. (2024).
Payment Card Industry Data Security Standard, Version 4.0.1.
Published 11 June 2024.

\bibitem{hellaswag}
Zellers, R., Holtzman, A., Bisk, Y., Farhadi, A., \& Choi, Y. (2019).
HellaSwag: Can a Machine Really Finish Your Sentence?
In \emph{Proceedings of the 57th Annual Meeting of the Association for
Computational Linguistics}, pages 4791--4800, Florence, Italy.
Association for Computational Linguistics. (arXiv:1905.07830)

\bibitem{iso27001}
International Organization for Standardization / International
Electrotechnical Commission. (2022).
ISO/IEC 27001:2022 --- Information security, cybersecurity and privacy
protection --- Information security management systems ---
Requirements.

\end{thebibliography}

\end{document}
